\documentclass[10pt,a4paper]{amsart}
\usepackage[margin=19mm]{geometry}
\usepackage{amsmath,amssymb,mathtools}
\usepackage[scr=rsfs,cal=euler]{mathalfa}
\usepackage{xfrac}
\usepackage[unicode,hidelinks,bookmarks=false]{hyperref}
\newcommand{\ie}{i.e.\ }
\newcommand{\cf}{cf.\ }
\newcommand{\resp}{respectively\ }
\newcommand{\Subsection}{\subsection}
\providecommand{\nobreakdash}{\nobreak\hbox{-}}
\setlength{\emergencystretch}{2em}
\multlinegap0pt
\newtheorem{thm}{Theorem}
\title{First-order Schwartian derivative and some classes of univalent functions}
\author{Milutin Obradovi\'{c}, Nikola Tuneski}
\date{Source: arXiv:2606.03433v1 (CC BY 4.0)}
\begin{document}
\maketitle

\noindent\emph{Source and licence.} First-order Schwartian derivative and some classes of univalent functions. Version arXiv:2606.03433v1 (CC BY 4.0), distributed by arXiv under the Creative Commons Attribution 4.0 International licence, http://creativecommons.org/licenses/by/4.0/, as recorded in the arXiv licence metadata for this exact version and shown on the abstract page https://arxiv.org/abs/2606.03433v1. Full licence notice: \texttt{licenses/arxiv-2606.03433-CC-BY-4.0.txt}.\par\smallskip

\noindent\textit{Editorial scope.} Counted unit: the article's thm th5 (source0.tex lines 445--447). The article's complete original proof of it is delivered with it (source0.tex lines 449--505, one \texttt{proof} environment of 57 lines). No mathematics is written, repaired or re-derived: the statement and the proof are the article's own lines, in the article's own notation, and no retained line is substituted or re-flowed.  Retained uncounted as the source-local context the proof needs: lines 158--160; lines 424--426; lines 432--439.  Nothing was omitted from any retained span, so the package carries no omission record.  No retained line contains a citation, so the wrapper carries no bibliography and no bibliography entry.  No retained line contains a figure environment, an image-inclusion command or drawing code, so no image is delivered and the package declares no asset dependency.  Editorial formatting only, in addition: an AMS \texttt{amsart} preamble that restates the article's own declarations and macros used by the retained lines, namely the environments \texttt{thm} on separate counters, together with the standard packages those lines need.  The retained environments print in this document as Theorem 1 in order of appearance, whereas the article prints the same items with the numbering of its own full document.  The complete native source of the article as distributed in the arXiv e-print for this exact version is bundled unchanged as \texttt{sources/201973/source0.tex}.
\begin{equation}\label{eq-2}
S'_{f}(0)=24(a_{4}-2a_{2}a_{3}+a_{2}^{3}).
\end{equation}

\begin{equation}\label{eq-16}
  |\omega_{11}|^2 + 3 |\omega_{13}|^2 + 5|\omega_{15}|^2 \leq1.
\end{equation}

\begin{equation}\label{eq-17}
\begin{split}
a_{2}&=2\omega _{11},\\
a_{3}&=2\omega_{13}+3\omega_{11}^{2}, \\
a_{4}&=2\omega_{33}+8\omega_{11}\omega_{13}+\frac{10}{3}\omega_{11}^{3},\\
0&=3\omega_{15}-3\omega_{11}\omega_{13}+\omega_{11}^{3}-3\omega_{33}.
\end{split}
\end{equation}

\begin{thm}\label{th5}
If $f$ belongs to $\mathcal{S}$, then $|S'_{f}(0)|\leq \frac{64}{5}\sqrt{3}=22.17025\ldots.$
\end{thm}

\begin{proof}
Using \eqref{eq-2} and \eqref{eq-17}, we have
\[
\begin{split}
|S'_{f}(0)|&=24|a_{4}-2a_{2}a_{3}+a_{2}^{3}|\\
&=48\left|\omega_{33}-\frac{1}{3}\omega_{11}^{3}\right|,
\end{split}
\]
or, if we use the fourth relation from \eqref{eq-17} and the relation \eqref{eq-16},
\[
\begin{split}
|S'_{f}(0)|&=48|\omega_{15}-\omega_{11}\omega_{13}|\\
&\leq48(|\omega_{15}|+|\omega_{11}||\omega_{13}|)\\
&\leq 48\left(\frac{1}{\sqrt{5}}\sqrt{1-|\omega_{11}|^2-3|\omega_{13}|^2}+|\omega_{11}||\omega_{13}|\right)\\
&=:48\cdot g(|\omega_{11}|,|\omega_{13}|),
\end{split}
\]
where
$g(x,y)=\frac{1}{\sqrt{5}}\sqrt{1-x^2-3y^2}+x y$ and
\[ (x,y)\in \Omega = \left\{(x,y) : 0\leq x \leq 1,\, 0\leq y\leq \frac{1}{\sqrt3}\sqrt{1-x^2} \right\}. \]

\medskip

Now, let find maximal value of the function $g$ on $\Omega$.

\medskip

To derive the critical points of $g$, we need to solve in the interior of $\Omega$ the following system
\[
\left\{
\begin{split}
\frac{\partial g(x,y)}{\partial x} &= y-\frac{1}{\sqrt{5} }\frac{x}{\sqrt{1-x^2-3 y^2}} = 0,\\
\frac{\partial g(x,y)}{\partial y} &= x-\frac{1}{\sqrt{5} }\frac{3 y}{\sqrt{1-x^2-3 y^2}} = 0.
\end{split}
\right.
\]
Now, from
\[3y\frac{\partial f_2(x,y)}{\partial x} - x \frac{\partial f_2(x,y)}{\partial y} =
3 y^2-x^2 = 0 , \]
we receive
\[y^2=\frac{x^2}{3}.\]
Substituting it in the second equation of the system gives
\[ \frac{x}{\sqrt{3}}-\frac{1}{\sqrt{5} }\frac{x}{\sqrt{1-2 x^2}} = 0,\]
which brings the unique solution of the system in the interior of $\Omega$, $x_1=\frac{1}{\sqrt5}$ and $y_1=\frac{1}{\sqrt{15}}$ such that $g(x_1,y_1)=\frac{64}{5}\sqrt{3}=22.17025\ldots$.

\medskip

On the edges of $\Omega$ we have:
\begin{itemize}
  \item[-] $g(0,y)=\frac{1}{\sqrt5}\sqrt{1-3 y^2}\le \frac{1}{\sqrt5} = 0.4472\ldots;$
  \item[-] $g(1,y) = g(1,0) = 0$;
  \item[-] $g(x,0) = \frac{1}{\sqrt5}\sqrt{1-x^2}\le \frac{1}{\sqrt5} = 0.4472\ldots;$
  \item[-] $h(x):=g\left(x, \sqrt{1-x^2}/\sqrt3\right) = \frac{1}{\sqrt3}x\sqrt{1-x^2}\le h(1/\sqrt2) = \sqrt3/6 = 0.288675\ldots$.
\end{itemize}
Combining the above analysis, we receive that the function $g$, on the domain $\Omega$ achieves the greatest value $\frac{64}{5}\sqrt{3}$ obtained for $x_1=\frac{1}{\sqrt5}$ and $y_1=\frac{1}{\sqrt{15}}$, i.e.,
\[|S'_{f}(0)|\leq \frac{64}{5}\sqrt{3}=22.17025\ldots.\]
\end{proof}

\end{document}
